\chapter{从研究问题到可部署系统}

本章讨论一个常被低估的事实：模型训练只是数据智能项目的一部分。工业问题通常包含观测、标签、时间、设备、成本和安全约束。一个在随机划分测试集上表现优秀的模型，如果不能在正确的时间点获得输入、不能解释告警来源，或无法承受缺失数据，仍然不是可用的工业系统。

\section{研究问题的形式化}
一个完整的问题定义至少应包含输入随机变量$X$、目标$Y$、决策$A$、环境状态$S$和损失$\ell$。预测任务可以写为
\begin{equation}
 f^*=\argmin_f\mathbb{E}_{(X,Y)\sim P_{deploy}}[\ell(f(X),Y)],
\end{equation}
而决策任务还需要考虑行动改变环境分布的反馈。训练数据来自$P_{train}$，部署数据来自$P_{deploy}$；二者不相同时，泛化问题就不再只是有限样本误差。

研究问题应写成可证伪的假设。例如，“加入振动频谱后，跨设备故障识别的宏平均 F1 在置信区间意义下提高”比“多模态模型效果更好”更适合作为研究假设。前者明确了输入、任务、评价量和比较范围。

\section{数据划分与泄漏}
随机划分适用于样本独立同分布且样本之间没有强关联的情形。工业数据常按设备、批次、工单或时间成组。若同一设备的相邻窗口同时出现在训练集和测试集，模型可能只需识别设备指纹即可获得高分。

设每条样本属于组$g$，则组级划分要求
\begin{equation}
 \mathcal{G}_{train}\cap\mathcal{G}_{test}=\varnothing.
\end{equation}
时间预测还必须满足训练样本时间早于测试样本时间。特征工程中的归一化统计量、缺失值填充规则和词表也只能由训练集估计，否则测试信息会通过预处理泄漏。

\section{基线的层次}
基线不是形式上的陪跑模型，而是解释改进来源的参照系。建议至少包含四层：
\begin{enumerate}[label=\arabic*.]
 \item 业务规则或历史平均，回答模型是否超过当前流程；
 \item 线性模型、树模型或最近邻，回答复杂模型是否必要；
 \item 文献中最强的可复现方法，回答实现是否具有研究竞争力；
 \item 与当前方法参数量、预训练数据和输入信息尽量匹配的强基线。
\end{enumerate}
如果新模型使用了额外传感器、更长时间窗口或更多预训练数据，就不能把性能提升全部归因于模型结构。

\section{消融实验的因果思路}
消融应一次只改变一个主要因素。若完整模型为$M(A,B,C)$，至少需要比较$M(A,B,0)$、$M(A,0,C)$和$M(0,B,C)$，并报告重复运行的均值与标准差。对于模块之间可能有交互的情形，应增加二因素实验，而不是只观察去掉全部模块后的巨大下降。

\section{不确定性与统计报告}
单次运行的准确率不能代表算法的稳定性能。设第$r$次独立运行得到指标$m_r$，样本均值和标准误为
\begin{equation}
 \bar m=\frac1R\sum_{r=1}^R m_r,
 \qquad
 \operatorname{SE}(\bar m)=\frac{s_m}{\sqrt R}.
\end{equation}
对于强相关的时间序列，不能简单把窗口数量当作独立样本数。更合理的做法是按设备或时间块进行重采样，并在测试集只使用一次的前提下报告置信区间。统计显著不等于工程显著，还需给出误报、漏报、延迟和成本变化。

\section{误差分析模板}
误差分析应回答三个问题：模型在哪些条件下出错，错误是否集中于某些群体或设备，错误能否通过可观测信息被提前识别。可按工况、设备年龄、缺失率、信号质量、故障类型和预测提前量分层。对于分类任务，混淆矩阵应配合样本量；对于回归任务，应画出残差随时间和工况的变化，而不是只报告一个平均误差。

\section{可复现训练协议}
一次可复现实验至少保存：代码版本、依赖版本、随机种子、数据快照或数据哈希、划分文件、配置文件、训练日志、最优检查点、评估脚本和硬件信息。随机种子并不能保证所有 GPU 运算完全确定，因此还应记录软件和硬件环境。对数据不能公开的项目，可以公开生成的索引规则、字段定义、脱敏统计量和最小可运行示例。

\begin{table}[htbp]
\centering
\caption{实验记录的最小字段}
\begin{tabular}{p{0.25\textwidth}p{0.62\textwidth}}
\toprule
字段 & 应记录的内容\\
\midrule
数据 & 时间范围、设备范围、标签规则、缺失率和划分方式\\
模型 & 参数量、输入窗口、初始化、冻结层和损失权重\\
优化 & 优化器、学习率、批大小、调度器、早停规则\\
评估 & 指标定义、阈值选择、置信区间和分层结果\\
系统 & 推理延迟、显存、吞吐、故障回退和监控指标\\
\bottomrule
\end{tabular}
\end{table}

\section{部署前的分布漂移检查}
部署前应建立训练分布画像，包括每个特征的分位数、缺失模式、类别频率和联合关系。上线后按时间监控这些统计量。对于特征分布，可使用分箱后的 PSI 或 Wasserstein 距离；对于标签延迟的任务，先监控输入漂移和模型置信度，再在标签到达后评估性能漂移。漂移报警不等于模型必然失效，但它应触发调查。

\section{阈值、成本与人工闭环}
分类器输出概率并不直接等价于行动。设假阳性成本为$C_{FP}$，假阴性成本为$C_{FN}$，告警阈值应依据
\begin{equation}
 \mathbb{E}[C\mid x]=C_{FP}P(\hat y=1,y=0\mid x)+C_{FN}P(\hat y=0,y=1\mid x)
\end{equation}
选择。阈值可以按设备、工况和维护资源动态调整，但每次调整都要保留审计记录。人工复核结果还可作为主动学习样本，不过必须防止只收集模型最容易发现的错误。

\section{研究阅读与复现路线}
阅读一篇论文时，先写出问题定义、数据生成过程、训练目标、比较对象和证据链，再阅读实现细节。复现顺序建议为：先运行作者提供的最小示例，再固定数据划分，随后复现主表，最后进行模块替换。若主结果无法复现，应区分数据版本、预处理、随机性、硬件和实现差异，不应立即修改方法以追逐一个数字。

\section{练习}
\begin{enumerate}[label=\arabic*.]
 \item 为预测性维护任务写出一个可证伪的研究假设。
 \item 给出一个会导致工业时间序列泄漏的预处理例子。
 \item 设计包含规则、线性模型、树模型和深度模型的基线表。
 \item 为漏报成本远高于误报成本的任务选择评价指标和阈值策略。
 \item 写出一个最小实验配置文件应包含的字段。
\end{enumerate}